\studySubsection{practice-calc-13}{练习库：高等数学第 13 讲 多元函数微分学}

\problemAnchor{MATH1-CALC-0081}
\begin{problemBox}
\begin{problemMeta}
\textbf{题目编号：} \problemRef{MATH1-CALC-0081} \quad
\textbf{答案：} \answerRef{MATH1-CALC-0081} \quad
\textbf{来源：} Codex 原创 / K132 深度讲义配套练习（非真题）
\end{problemMeta}
\begin{enumerate}
\item[\textbf{A}] 求 \(f(x,y)=\sqrt{1-x^2-4y^2}\) 的定义域，并说明边界是否包含。
\item[\textbf{B}] 求
\[
f(x,y)=\frac{\ln(y-x)}{\sqrt{\,9-x^2-y^2\,}}
\]
的定义域。
\item[\textbf{C}] 说明 \(z=4-x^2-y^2\) 的定义域、水平截线及图形类型。
\item[\textbf{M}] 设
\[
D_a=\{(x,y):x^2+y^2\le a,\ x+y>1\}.
\]
求实数 \(a\) 使 \(D_a\ne\varnothing\)。
\end{enumerate}
\end{problemBox}

\problemAnchor{MATH1-CALC-0082}
\begin{problemBox}
\begin{problemMeta}
\textbf{题目编号：} \problemRef{MATH1-CALC-0082} \quad
\textbf{答案：} \answerRef{MATH1-CALC-0082} \quad
\textbf{来源：} Codex 原创 / K133 深度讲义配套练习（非真题）
\end{problemMeta}
\begin{enumerate}
\item[\textbf{A}] 求 \(\displaystyle\lim_{(x,y)\to(0,0)}(2x^2+3y^2)\)。
\item[\textbf{B}] 用统一估计求
\[
\lim_{(x,y)\to(0,0)}\frac{x^2y}{\sqrt{x^2+y^2}}.
\]
\item[\textbf{C}] 证明
\[
\lim_{(x,y)\to(0,0)}
\frac{x^4+y^4}{\sqrt{x^2+y^2}}=0.
\]
\item[\textbf{M}] 直接按 \(\varepsilon\)-\(\delta\) 定义证明
\(\lim_{(x,y)\to(0,0)}(x^2+3y^2)=0\)，并给出一个可用的
\(\delta(\varepsilon)\)。
\end{enumerate}
\end{problemBox}

\problemAnchor{MATH1-CALC-0083}
\begin{problemBox}
\begin{problemMeta}
\textbf{题目编号：} \problemRef{MATH1-CALC-0083} \quad
\textbf{答案：} \answerRef{MATH1-CALC-0083} \quad
\textbf{来源：} Codex 原创 / K134 深度讲义配套练习（非真题）
\end{problemMeta}
\begin{enumerate}
\item[\textbf{A}] 判断
\[
\lim_{(x,y)\to(0,0)}
\frac{x^2+xy+y^2}{x^2+y^2}
\]
是否存在。
\item[\textbf{B}] 判断
\[
\lim_{(x,y)\to(0,0)}
\frac{x^2\bigl(y-\sin x\bigr)}
{x^4+\bigl(y-\sin x\bigr)^2}
\]
是否存在，并说明如何从式子结构选择曲线路径。
\item[\textbf{C}] 在定义域
\[
D=\{(x,y):y\ge |x|\}\setminus\{(0,0)\}
\]
内，判断
\[
\lim_{(x,y)\to(0,0)}\frac{x}{y}
\]
是否存在。
\item[\textbf{M}] 证明
\[
\sin\!\frac1{x^2+y^2}
\]
在 \((x,y)\to(0,0)\) 时不存在二重极限。要求给出两列趋于原点的点，
不得只写“函数振荡”。
\end{enumerate}
\end{problemBox}

\problemAnchor{MATH1-CALC-0084}
\begin{problemBox}
\begin{problemMeta}
\textbf{题目编号：} \problemRef{MATH1-CALC-0084} \quad
\textbf{答案：} \answerRef{MATH1-CALC-0084} \quad
\textbf{来源：} Codex 原创 / K135 深度讲义配套练习（非真题）
\end{problemMeta}
\begin{enumerate}
\item[\textbf{A}] 求
\(\displaystyle\lim_{(x,y)\to(0,0)}x^3/(x^2+y^2)\)。
\item[\textbf{B}] 求
\[
\lim_{(x,y)\to(0,0)}
\frac{x^4y^2}{(x^2+y^2)^2}.
\]
\item[\textbf{C}] 求
\[
\lim_{(x,y)\to(0,0)}
\frac{\sin(x^2+y^2)}{\sqrt{x^2+y^2}}.
\]
\item[\textbf{M}] 对 \(\alpha\in\mathbb R\)，讨论
\[
f_\alpha(x,y)
=(x^2+y^2)^{\alpha/2}\frac{x^2-y^2}{x^2+y^2}
\]
在 \((x,y)\to(0,0)\) 时的极限。
\end{enumerate}
\end{problemBox}

\problemAnchor{MATH1-CALC-0085}
\begin{problemBox}
\begin{problemMeta}
\textbf{题目编号：} \problemRef{MATH1-CALC-0085} \quad
\textbf{答案：} \answerRef{MATH1-CALC-0085} \quad
\textbf{来源：} Codex 原创 / K136 深度讲义配套练习（非真题）
\end{problemMeta}
\begin{enumerate}
\item[\textbf{A}] 判断 \(f(x,y)=e^{x-y}+x^2y\) 在 \(\mathbb R^2\) 上的连续性。
\item[\textbf{B}] 令
\[
f(x,y)=
\begin{cases}
\dfrac{xy^2}{x^2+y^2},&(x,y)\ne(0,0),\\
0,&(x,y)=(0,0).
\end{cases}
\]
判断 \(f\) 在原点是否连续。
\item[\textbf{C}] 只利用存在性定理证明
\(f(x,y)=x^2+2xy+3y^2\) 在三角形
\(D=\{x\ge0,y\ge0,x+y\le1\}\) 上取得最大值和最小值。
\item[\textbf{M}] 分别给出一个反例，说明最值存在性定理中“区域闭”和“函数连续”
任一条件删除后，最大值都可能不存在。第一例不得使用圆盘。
\end{enumerate}
\end{problemBox}

\problemAnchor{MATH1-CALC-0086}
\begin{problemBox}
\begin{problemMeta}
\textbf{题目编号：} \problemRef{MATH1-CALC-0086} \quad
\textbf{答案：} \answerRef{MATH1-CALC-0086} \quad
\textbf{来源：} Codex 原创 / K137 深度讲义配套练习（非真题）
\end{problemMeta}
\begin{enumerate}
\item[\textbf{A}] 求 \(f(x,y)=x^3y+\cos(xy)\) 的 \(f_x,f_y\)。
\item[\textbf{B}] 按定义求 \(f(x,y)=x|y|+y^2\) 在原点的两个偏导。
\item[\textbf{C}] 令
\[
f(x,y)=
\begin{cases}
\dfrac{x^2-y^2}{x^2+y^2},&(x,y)\ne(0,0),\\
0,&(0,0).
\end{cases}
\]
判断原点两个偏导是否存在，并判断函数是否连续。
\item[\textbf{M}] 对 \(\alpha>0\)，讨论
\(f_\alpha(x,y)=|x|^\alpha+y^2\) 在原点的 \(x\) 偏导是否存在。
\end{enumerate}
\end{problemBox}

\problemAnchor{MATH1-CALC-0087}
\begin{problemBox}
\begin{problemMeta}
\textbf{题目编号：} \problemRef{MATH1-CALC-0087} \quad
\textbf{答案：} \answerRef{MATH1-CALC-0087} \quad
\textbf{来源：} Codex 原创 / K138 深度讲义配套练习（非真题）
\end{problemMeta}
\begin{enumerate}
\item[\textbf{A}] 求 \(f(x,y)=\sin(xy)\) 的四个二阶偏导。
\item[\textbf{B}] 对 \(f=x^2y^3+e^x y\)，分别从两种次序计算 \(f_{xy},f_{yx}\)。
\item[\textbf{C}] 设
\[
f(x,y)=
\begin{cases}
\dfrac{xy^3}{x^2+y^2},&(x,y)\ne(0,0),\\0,&(0,0),
\end{cases}
\]
按定义计算 \(f_{xy}(0,0)\) 与 \(f_{yx}(0,0)\)。
\item[\textbf{M}] 已知某个 \(C^2\) 函数满足
\[
f_x=ay+xy^2,\qquad f_y=2x+x^2y.
\]
求参数 \(a\)，并写出一个这样的 \(f\)。
\end{enumerate}
\end{problemBox}

\problemAnchor{MATH1-CALC-0088}
\begin{problemBox}
\begin{problemMeta}
\textbf{题目编号：} \problemRef{MATH1-CALC-0088} \quad
\textbf{答案：} \answerRef{MATH1-CALC-0088} \quad
\textbf{来源：} Codex 原创 / K139 深度讲义配套练习（非真题）
\end{problemMeta}
\begin{enumerate}
\item[\textbf{A}] 求 \(z=x^2y+y^3\) 在点 \((1,-1)\) 的全微分。
\item[\textbf{B}] 用线性化估算
\(\sqrt{9-(1.02)^2-(1.01)^2}\)，并写明线性化点。
\item[\textbf{C}] 从定义证明 \(f(x,y)=x^2+xy+y^2\) 在原点可微并求全微分。
\item[\textbf{M}] 圆柱体体积 \(V=\pi r^2h\)。当 \(r,h\) 的小测量误差分别为
\(dr,dh\) 时，写出体积的一阶误差，并说明精确增量中被舍去项的阶。
\end{enumerate}
\end{problemBox}

\problemAnchor{MATH1-CALC-0089}
\begin{problemBox}
\begin{problemMeta}
\textbf{题目编号：} \problemRef{MATH1-CALC-0089} \quad
\textbf{答案：} \answerRef{MATH1-CALC-0089} \quad
\textbf{来源：} Codex 原创 / K140 深度讲义配套练习（非真题）
\end{problemMeta}
\begin{enumerate}
\item[\textbf{A}] 判断并改正：偏导存在 \(\Rightarrow\) 连续；
偏导连续 \(\Rightarrow\) 可微；可微 \(\Rightarrow\) 偏导存在。
\item[\textbf{B}] 验证
\[
f(x,y)=
\begin{cases}
\dfrac{x^2y^2}{x^4+y^4},&(x,y)\ne(0,0),\\0,&(0,0)
\end{cases}
\]
在原点两个偏导存在但不连续。
\item[\textbf{C}] 判断
\[
f(x,y)=
\begin{cases}
\dfrac{y^3}{x^2+y^2},&(x,y)\ne(0,0),\\0,&(0,0)
\end{cases}
\]
在原点是否可微。
\item[\textbf{M}] 证明
\[
f(x,y)=
\begin{cases}
x^2\sin\!\dfrac1x,&x\ne0,\\0,&x=0
\end{cases}
\]
在原点可微，但至少一个一阶偏导在原点不连续。
\end{enumerate}
\end{problemBox}

\problemAnchor{MATH1-CALC-0090}
\begin{problemBox}
\begin{problemMeta}
\textbf{题目编号：} \problemRef{MATH1-CALC-0090} \quad
\textbf{答案：} \answerRef{MATH1-CALC-0090} \quad
\textbf{来源：} Codex 原创 / K141 深度讲义配套练习（非真题）
\end{problemMeta}
\begin{enumerate}
\item[\textbf{A}] 按定义判断 \(f(x,y)=(x^2+y^2)^{2/3}\) 在原点是否可微。
\item[\textbf{B}] 判断
\[
f(x,y)=
\begin{cases}
\dfrac{|xy|}{\sqrt{x^2+y^2}},&(x,y)\ne(0,0),\\0,&(0,0)
\end{cases}
\]
在原点是否可微。
\item[\textbf{C}] 证明
\[
f(x,y)=3x+2y+(x^2+y^2)^{3/2}\cos\!\frac1{\sqrt{x^2+y^2}}
\]
（原点处余项定义为 \(0\)）在原点可微，并求全微分。
\item[\textbf{M}] 对 \(\alpha>0\)，讨论
\[
f_\alpha(x,y)=(x^2+y^4)^\alpha
\]
在原点的可微性。
\end{enumerate}
\end{problemBox}

\problemAnchor{MATH1-CALC-0091}
\begin{problemBox}
\begin{problemMeta}
\textbf{题目编号：} \problemRef{MATH1-CALC-0091} \quad
\textbf{答案：} \answerRef{MATH1-CALC-0091} \quad
\textbf{来源：} Codex 原创 / K142 深度讲义配套练习（非真题）
\end{problemMeta}
\begin{enumerate}
\item[\textbf{A}] 设 \(z=f(x-y,x^2+y^2)\)，求 \(z_x,z_y\)。
\item[\textbf{B}] 设 \(w=f(xy,yz,zx)\)，求 \(w_y\)。
\item[\textbf{C}] 设 \(z=f(e^x\cos y,e^x\sin y)\)，用 \(f_u,f_v\) 表示
\(z_x,z_y\)。
\item[\textbf{M}] 设 \(z=f(x+y,x-y)\)。证明
\[
z_x^2-z_y^2=4f_u f_v,
\]
并写清所有偏导的求值点。
\end{enumerate}
\end{problemBox}

\problemAnchor{MATH1-CALC-0092}
\begin{problemBox}
\begin{problemMeta}
\textbf{题目编号：} \problemRef{MATH1-CALC-0092} \quad
\textbf{答案：} \answerRef{MATH1-CALC-0092} \quad
\textbf{来源：} Codex 原创 / K143 深度讲义配套练习（非真题）
\end{problemMeta}
\begin{enumerate}
\item[\textbf{A}] 设 \(z=f(x^2+2y^2)\)，求 \(z_{xx},z_{yy}\)。
\item[\textbf{B}] 设 \(f\in C^2\)，且
\(z=f(2x+y,x-2y)\)，求 \(z_{xx}\)。
\item[\textbf{C}] 设 \(z=e^u,u=xy\)，分别由链式法则和直接求导计算 \(z_{xy}\)。
\item[\textbf{M}] 设 \(z=f(u),u=x^2+x+y^2\)。说明为什么
\(z_{xx}\ne f''(u)u_x^2\) 一般不成立，并写出正确公式。
\end{enumerate}
\end{problemBox}

\problemAnchor{MATH1-CALC-0093}
\begin{problemBox}
\begin{problemMeta}
\textbf{题目编号：} \problemRef{MATH1-CALC-0093} \quad
\textbf{答案：} \answerRef{MATH1-CALC-0093} \quad
\textbf{来源：} Codex 原创 / K144 深度讲义配套练习（非真题）
\end{problemMeta}
\begin{enumerate}
\item[\textbf{A}] 设 \(z=f(u,v),u=2x-y,v=x+3y\)，用 \(dx,dy\) 表示 \(dz\)。
\item[\textbf{B}] 对 \(z=u v^2,u=x+y,v=x-y\)，分别用形式不变性和直接展开求 \(dz\)。
\item[\textbf{C}] 在极坐标换元下求 \(z=x^2-y^2\) 的 \(dz\)，并化简为
\(dr,d\theta\) 的表达。
\item[\textbf{M}] 设
\[
z=u^2+v,\qquad u=x^2,\qquad v=y^2.
\]
分别直接对 \(z=x^4+y^2\) 求 \(d^2z\)，以及把一阶复合关系完整微分，
说明只代入外层二阶形式会漏掉哪些项。
\end{enumerate}
\end{problemBox}

\problemAnchor{MATH1-CALC-0094}
\begin{problemBox}
\begin{problemMeta}
\textbf{题目编号：} \problemRef{MATH1-CALC-0094} \quad
\textbf{答案：} \answerRef{MATH1-CALC-0094} \quad
\textbf{来源：} Codex 原创 / K145 深度讲义配套练习（非真题）
\end{problemMeta}
\begin{enumerate}
\item[\textbf{A}] 方程 \(x^2+2y^2+z^2=4\) 在点 \((1,1,1)\) 附近确定
\(z=z(x,y)\)。求 \(z_x,z_y\)。
\item[\textbf{B}] 由 \(xz+e^z+y^2=0\) 局部确定 \(z(x,y)\)。在
\(e^z+x\ne0\) 时求两个一阶偏导。
\item[\textbf{C}] 对方程 \(x^2+y^2+xyz=3\) 求 \(dz\)，并写明能解出
\(dz\) 的条件。
\item[\textbf{M}] 讨论 \(z^3=x^2+y^2\) 在原点为何不能直接使用隐函数偏导公式，
并判断实际分支是否唯一、是否可微。
\end{enumerate}
\end{problemBox}

\problemAnchor{MATH1-CALC-0095}
\begin{problemBox}
\begin{problemMeta}
\textbf{题目编号：} \problemRef{MATH1-CALC-0095} \quad
\textbf{答案：} \answerRef{MATH1-CALC-0095} \quad
\textbf{来源：} Codex 原创 / K146 深度讲义配套练习（非真题）
\end{problemMeta}
\begin{enumerate}
\item[\textbf{A}] 由 \(u+2v=x,\ 3u-v=x^2\) 确定 \(u(x),v(x)\)，求 \(u',v'\)。
\item[\textbf{B}] 方程组
\[
u^2+v=x,\qquad u+v^2=e^x
\]
局部确定 \(u,v\)。写出求 \(u',v'\) 的线性方程组及 Jacobian 条件。
\item[\textbf{C}] 方程组
\[
u+v+xy=0,\qquad u-v+x^2-y=0
\]
确定 \(u(x,y),v(x,y)\)。用隐式线性方程组求
\((u_x,v_x)\) 与 \((u_y,v_y)\)，并用显式解核验。
\item[\textbf{M}] 用 \(u^3=x,\ v=y\) 说明 Jacobian 在原点退化时，
即使实分支唯一，局部可微性仍可能失败。
\end{enumerate}
\end{problemBox}

\problemAnchor{MATH1-CALC-0096}
\begin{problemBox}
\begin{problemMeta}
\textbf{题目编号：} \problemRef{MATH1-CALC-0096} \quad
\textbf{答案：} \answerRef{MATH1-CALC-0096} \quad
\textbf{来源：} Codex 原创 / K147 深度讲义配套练习（非真题）
\end{problemMeta}
\begin{enumerate}
\item[\textbf{A}] 判断 \(f(x,y)=x^4+x^2y^2+y^4\) 是否齐次，并给出次数。
\item[\textbf{B}] 不完全展开偏导，求
\(x f_x+y f_y\)，其中 \(f=(x^2+y^2)^2\)。
\item[\textbf{C}] 设 \(a,b>0\)，且可微函数满足加权缩放关系
\[
f(t^a x,t^b y)=t^m f(x,y)\qquad(t>0).
\]
从定义推导相应的加权 Euler 公式。
\item[\textbf{M}] 若 \(f\in C^1\) 为 \(m\) 次齐次函数，证明
\(f_x,f_y\)（在其存在处）都是 \(m-1\) 次齐次函数。
\end{enumerate}
\end{problemBox}

\problemAnchor{MATH1-CALC-0097}
\begin{problemBox}
\begin{problemMeta}
\textbf{题目编号：} \problemRef{MATH1-CALC-0097} \quad
\textbf{答案：} \answerRef{MATH1-CALC-0097} \quad
\textbf{来源：} Codex 原创 / K148 深度讲义配套练习（非真题）
\end{problemMeta}
\begin{enumerate}
\item[\textbf{A}] 按定义求 \(f=x^2+2y^2\) 在 \((1,1)\) 沿方向 \((1,2)\) 的方向导数。
\item[\textbf{B}] 求 \(f=xy\) 在 \((2,-1)\) 沿与 \(x\) 轴夹角 \(\pi/3\) 的方向导数。
\item[\textbf{C}] 按双侧定义判断 \(f(x,y)=\sqrt{x^2+4y^2}\) 在原点沿任意单位方向的
方向导数是否存在。
\item[\textbf{M}] 令
\[
f(x,y)=
\begin{cases}\dfrac{x^2y}{x^4+y^2},&(x,y)\ne(0,0),\\0,&(0,0).\end{cases}
\]
证明原点各方向导数均存在，但 \(f\) 在原点不可微。
\end{enumerate}
\end{problemBox}

\problemAnchor{MATH1-CALC-0098}
\begin{problemBox}
\begin{problemMeta}
\textbf{题目编号：} \problemRef{MATH1-CALC-0098} \quad
\textbf{答案：} \answerRef{MATH1-CALC-0098} \quad
\textbf{来源：} Codex 原创 / K149 深度讲义配套练习（非真题）
\end{problemMeta}
\begin{enumerate}
\item[\textbf{A}] 求 \(f=x^2+2xy+3y^2\) 在 \((1,-1)\) 的梯度、最大方向导数及方向。
\item[\textbf{B}] 求 \(T=3x^2+y^2\) 在 \((1,2)\) 的最速下降方向和最小方向导数。
\item[\textbf{C}] 设
\[
f(x,y)=(y-x^2)(y-2x^2).
\]
说明原点梯度为零，却不是极值点。
\item[\textbf{M}] 求椭圆 \(x^2+2y^2=3\) 在点 \((1,1)\) 的法向量和切线方程。
\end{enumerate}
\end{problemBox}

\problemAnchor{MATH1-CALC-0099}
\begin{problemBox}
\begin{problemMeta}
\textbf{题目编号：} \problemRef{MATH1-CALC-0099} \quad
\textbf{答案：} \answerRef{MATH1-CALC-0099} \quad
\textbf{来源：} Codex 原创 / K150 深度讲义配套练习（非真题）
\end{problemMeta}
\begin{enumerate}
\item[\textbf{A}] 求椭球面 \(x^2+2y^2+z^2=4\) 在 \((1,1,1)\) 的切平面与法线。
\item[\textbf{B}] 求曲面 \(z=x^2-2y^2\) 在 \((1,1,-1)\) 的切平面。
\item[\textbf{C}] 在曲面 \(z=x^2+y^2\) 上求切平面平行于
\(4x+2y-z=0\) 的点。
\item[\textbf{M}] 对曲面
\[
z^2=x^4+y^4
\]
说明原点处梯度公式为何失效，并判断该曲面是否仍有唯一切平面。
用两支显函数的可微性作几何核验。
\end{enumerate}
\end{problemBox}
